% --------------------------------------------------------------
% Paquetes, configuración del documento, ...
% --------------------------------------------------------------
 
\documentclass[12pt]{article}
 
\usepackage[margin=1in]{geometry} 
\usepackage{amsmath,amsthm,amssymb}
\usepackage[spanish]{babel}
\addto\captionsspanish{\renewcommand{\tablename}{Tabla}}
\usepackage{graphicx}
\usepackage{subcaption}
\usepackage{enumitem}
\usepackage{hyperref}
\usepackage{listings}
\addto\captionsspanish{\renewcommand{\lstlistingname}{Listado}}
\usepackage{xcolor}
\usepackage{booktabs}

\lstset{
    language=C,
    basicstyle=\ttfamily\small,
    keywordstyle=\color{blue!70!black}\bfseries,
    commentstyle=\color{gray}\itshape,
    stringstyle=\color{orange!80!black},
    backgroundcolor=\color{gray!8},
    frame=single,
    framerule=0.4pt,
    rulecolor=\color{gray!40},
    breaklines=true,
    captionpos=b,
    xleftmargin=8pt,
    xrightmargin=8pt,
    morekeywords={__global__, blockIdx, blockDim, threadIdx},
}
 
% \newcommand{\N}{\mathbb{N}}
% \newcommand{\Z}{\mathbb{Z}}
 
\newenvironment{exercise}[2][Ejercicio]{\begin{trivlist}
\item[\hskip \labelsep {\bfseries #1}\hskip \labelsep {\bfseries #2.}]}{\end{trivlist}}
 
\begin{document}

% --------------------------------------------------------------
%                  Zona de ejercicios/trabajo
% --------------------------------------------------------------

\title{Práctica 3 - Organización de \textit{threads} y patrones de acceso a memoria}
\author{Álvaro Martínez Alfaro\\
    Aceleradores Gráficos}

\maketitle

\begin{exercise}{1}
    Implementación paralela con CUDA de la suma de dos vectores de $n$ números enteros, usando la
    memoria global de la GPU. Se proporciona un programa secuencial de la suma de vectores. El
    acceso a los vectores se producirá de acuerdo a dos patrones diferentes:
    \begin{enumerate}[label=\alph*)]
        \item Cada thread procesa $m$ componentes consecutivas de los vectores.
        \item Cada thread procesa $m$ componentes de los vectores separadas una determinada
              distancia.
    \end{enumerate}
    El alumno debe obtener dos kernels, uno para cada caso, y probar su correcto funcionamiento.
    Para esta primera parte se puede usar cualquier valor de $n$ y $m$, aunque es recomendable
    probar con varios. Para tamaños grandes de los vectores, se comprobará la corrección de los
    kernels comparando el vector suma resultado con el obtenido con la versión secuencial.

    \begin{center}
        \noindent\rule{4cm}{0.2pt}
    \end{center}

    Se han implementado dos kernels CUDA, \texttt{SumVectorA} y \texttt{SumVectorB}, ambos
    calculando $C = A + B$ sobre vectores de $n$ enteros almacenados en memoria global de la GPU.
    La configuración de ejecución es común a los dos kernels: cada bloque contiene \texttt{block\_size}
    threads, y el número de bloques se calcula como $\lceil n\,/\,(\texttt{block\_size} \cdot m) \rceil$,
    de forma que el total de threads cubre exactamente los $n$ elementos repartiendo $m$ elementos
    por thread.

    \subsubsection*{Kernel \texttt{SumVectorA}: acceso a componentes consecutivas}

    Para el kernel \texttt{SumVectorA}, cada thread tiene asignado un rango de $m$ posiciones
    contiguas del vector. El índice de inicio se calcula a partir del identificador global del thread:
    %
    \begin{equation}
        i = (\texttt{blockIdx.x} \cdot \texttt{blockDim.x} + \texttt{threadIdx.x}) \cdot m
    \end{equation}
    %
    El thread accede así a los elementos $i,\, i+1,\, \ldots,\, i+m-1$. La Tabla~\ref{tab:kernelA}
    muestra el reparto para $m = 4$ y \texttt{block\_size} $= 8$.

    \begin{table}[h]
        \centering
        \begin{tabular}{lll}
            \hline
            \textbf{Bloque} & \textbf{Thread} & \textbf{Elementos} \\
            \hline
            0               & 0               & $0,1,2,3$          \\
            0               & 1               & $4,5,6,7$          \\
            0               & $\vdots$        & $\vdots$           \\
            0               & 7               & $28,29,30,31$      \\
            1               & 0               & $32,33,34,35$      \\
            $\vdots$        & $\vdots$        & $\vdots$           \\
            \hline
        \end{tabular}
        \caption{Reparto de elementos en el Kernel A con $m=4$, \texttt{block\_size}$=8$.}
        \label{tab:kernelA}
    \end{table}

    \begin{lstlisting}[caption={Kernel \texttt{SumVectorA}: acceso consecutivo.}, label={lst:SumVectorA}]
__global__ void SumVectorA(vector A, vector B, vector C,
                            unsigned int n, int comp)
{
    int i = (blockIdx.x * blockDim.x + threadIdx.x) * comp;
    for (int j = 0; j < comp && i + j < n; j++)
        C[i + j] = A[i + j] + B[i + j];
}
    \end{lstlisting}

    Como se puede observar en el Listado \ref{lst:SumVectorA}, la condición \texttt{i+j<n} garantiza
    que el último thread no accede fuera de los límites del vector cuando $n$ no es múltiplo de
    $m \cdot \texttt{block\_size}$.

    \subsubsection*{Kernel \texttt{SumVectorB}: acceso a componentes separadas una distancia fija}

    En este kernel cada thread procesa $m$ elementos separados por una distancia igual al número
    de threads por bloque (\texttt{dist = block\_size}). El índice base del thread dentro del
    segmento de su bloque es:
    %
    \begin{equation}
        i = \texttt{blockIdx.x} \cdot \texttt{blockDim.x} \cdot m + \texttt{threadIdx.x}
    \end{equation}
    %
    y los elementos accedidos son $i,\, i+\texttt{dist},\, i+2\cdot\texttt{dist},\,\ldots,\,
        i+(m-1)\cdot\texttt{dist}$.
    La Tabla~\ref{tab:kernelB} muestra el reparto para $m = 4$ y \texttt{block\_size} $= 4$.

    \begin{table}[h]
        \centering
        \begin{tabular}{lll}
            \hline
            \textbf{Bloque} & \textbf{Thread} & \textbf{Elementos} \\
            \hline
            0               & 0               & $0,4,8,12$         \\
            0               & 1               & $1,5,9,13$         \\
            0               & 2               & $2,6,10,14$        \\
            0               & 3               & $3,7,11,15$        \\
            1               & 0               & $16,20,24,28$      \\
            $\vdots$        & $\vdots$        & $\vdots$           \\
            \hline
        \end{tabular}
        \caption{Reparto de elementos en el Kernel B con $m=4$, \texttt{block\_size}$=4$.}
        \label{tab:kernelB}
    \end{table}

    \begin{lstlisting}[caption={Kernel \texttt{SumVectorB}: acceso con saltos.}, label={lst:SumVectorB}]
__global__ void SumVectorB(vector A, vector B, vector C,
                            unsigned int n, int comp, int dist)
{
    int i = blockIdx.x * blockDim.x * comp + threadIdx.x;
    if (i >= n) return;
    for (int j = 0; j < comp; j++) {
        int idx = i + j * dist;
        if (idx < n)
            C[idx] = A[idx] + B[idx];
    }
}
    \end{lstlisting}

    Como se observa en el Listado \ref{lst:SumVectorB}, dentro de cada bloque, los threads acceden a
    posiciones contiguas en cada iteración $j$ (threads $0, 1, \ldots, \texttt{block\_size}-1$
    acceden a $i, i+1, \ldots$ con $i$ fijo para la iteración), lo que tiene implicaciones sobre la
    coalescencia que se analizan en el Ejercicio~2.

    \subsubsection*{Verificación de la corrección}

    Se verificó el funcionamiento de ambos kernels con distintos valores de $n$ y $m$. Para
    vectores pequeños se usó el modo de depuración (\texttt{debug=1}), que imprime los vectores
    de entrada y el resultado, permitiendo comprobar manualmente que $C[i] = A[i] + B[i]$.
    Para vectores grandes se comparó el resultado del kernel con el de la versión secuencial,
    obteniendo resultados idénticos en todos los casos probados.

    A modo de ejemplo, la Tabla~\ref{tab:verificacion} muestra una ejecución con $n=32$, $m=4$,
    \texttt{block\_size}$=8$, kernel A.

    \begin{table}[h]
        \centering
        \begin{tabular}{l *{11}{r}}
            \toprule
                    & \multicolumn{11}{c}{\textit{Índice}}                                                                  \\
            \cmidrule(l){2-12}
                    & $0$                                  & $1$ & $2$ & $3$ & $4$ & $5$ & $6$ & $7$ & $8$ & $9$ & $\cdots$ \\
            \midrule
            $A$     & 6                                    & 8   & 2   & 1   & 6   & 0   & 5   & 4   & 1   & 0   & $\cdots$ \\
            $B$     & 7                                    & 3   & 5   & 4   & 3   & 1   & 0   & 4   & 1   & 2   & $\cdots$ \\
            $C=A+B$ & 13                                   & 11  & 7   & 5   & 9   & 1   & 5   & 8   & 2   & 2   & $\cdots$ \\
            \bottomrule
        \end{tabular}
        \caption{Ejemplo de verificación: $n=32$, $m=4$, \texttt{block\_size}$=8$, kernel A.}
        \label{tab:verificacion}
    \end{table}

\end{exercise}

\begin{exercise}{2}
    Análisis estático de los dos kernels (sin realizar pruebas con ellos), desde el punto de vista
    del patrón de acceso a memoria. Para ello se debe estudiar la coalescencia y la localidad, y
    razonar sobre cuál de los kernels tendrá mejor rendimiento. El análisis se debe hacer teniendo
    en cuenta el valor de $m$, considerando diferentes valores de $m$ (desde pequeños a grandes, por
    ejemplo).

    \begin{center}
        \noindent\rule{4cm}{0.2pt}
    \end{center}

    En arquitecturas CUDA, un \textit{warp} agrupa 32 threads que ejecutan la misma instrucción
    simultáneamente. Cuando esos 32 threads acceden a posiciones de memoria contiguas y alineadas,
    la GPU los atiende con una única transacción de memoria (\textit{acceso coalescente}). Si en
    cambio los accesos están separados por un salto (\textit{stride}), se necesitan múltiples
    transacciones, lo que penaliza el rendimiento.

    \subsubsection*{Kernel \texttt{SumVectorA}: coalescencia y localidad}

    Cada thread tiene un identificador global $t = \texttt{blockIdx.x} \cdot \texttt{blockDim.x}
        + \texttt{threadIdx.x}$. En la iteración $j$ del bucle interno, el thread $t$ accede al índice:
    %
    \begin{equation}
        \text{idx}_A(t,j) = t \cdot m + j
    \end{equation}
    %
    Dos threads consecutivos del mismo warp, $t$ y $t+1$, acceden en la misma iteración $j$ a los
    índices $t \cdot m + j$ y $(t+1) \cdot m + j$, cuya distancia es $m$. Por tanto, el
    \textbf{stride entre threads es $m$}:

    \begin{itemize}
        \item \textbf{$m = 1$}: los 32 threads del warp acceden a posiciones $t, t+1, \ldots, t+31$,
              completamente contiguas. El acceso es \textbf{coalescente}.
        \item \textbf{$m > 1$}: la distancia entre accesos es $m$ elementos, de modo que el warp
              necesita hasta $\min(m, 32)$ transacciones de memoria por iteración en lugar de una.
              El acceso es \textbf{no coalescente}.
    \end{itemize}

    En cuanto a la localidad, cada thread individual accede a $m$ posiciones contiguas, lo que
    proporciona buena \textit{localidad espacial por thread}; sin embargo, a nivel de warp este
    beneficio desaparece para $m > 1$ por el problema de coalescencia descrito.

    \subsubsection*{Kernel \texttt{SumVectorB}: coalescencia y localidad}

    El índice base del thread con identificador local $\texttt{threadIdx.x} = t$ en el bloque
    $b = \texttt{blockIdx.x}$ es $i = b \cdot \texttt{blockDim.x} \cdot m + t$. En la iteración
    $j$, dicho thread accede al índice:
    %
    \begin{equation}
        \text{idx}_B(t,j) = b \cdot \texttt{blockDim.x} \cdot m + t + j \cdot \texttt{blockDim.x}
    \end{equation}
    %
    Dos threads consecutivos del warp, $t$ y $t+1$, acceden en la misma iteración $j$ a los índices
    $(\ldots + t + j \cdot \texttt{blockDim.x})$ y $(\ldots + t+1 + j \cdot \texttt{blockDim.x})$,
    cuya distancia es siempre $1$. Por tanto, \textbf{el acceso es coalescente para cualquier valor
        de $m$}.

    La contrapartida es que cada thread individual accede a posiciones separadas por el valor de
    \texttt{blockDim.x} elementos, lo que implica \textit{escasa localidad espacial por thread}.
    Sin embargo, dado que la coalescencia a nivel de warp está garantizada, este efecto no introduce
    penalización adicional en el ancho de banda de memoria.

    \subsubsection*{Predicción de rendimiento}

    La Tabla~\ref{tab:coalescencia} resume el análisis para los valores de $m$ considerados en el
    Ejercicio~3; esto es, $m \in \{1, 64, 256\}$.

    \begin{table}[h]
        \centering
        \begin{tabular}{ccccc}
            \toprule
            $m$   & \multicolumn{2}{c}{\textbf{Kernel A}} & \multicolumn{2}{c}{\textbf{Kernel B}}                                \\
            \cmidrule(lr){2-3}\cmidrule(lr){4-5}
                  & Stride en warp                        & Coalescente                           & Stride en warp & Coalescente \\
            \midrule
            $1$   & $1$                                   & Sí                                    & $1$            & Sí          \\
            $64$  & $64$                                  & No                                    & $1$            & Sí          \\
            $256$ & $256$                                 & No                                    & $1$            & Sí          \\
            \bottomrule
        \end{tabular}
        \caption{Coalescencia de los accesos a memoria en función de $m$.}
        \label{tab:coalescencia}
    \end{table}

    Para $m = 1$ ambos kernels son equivalentes: los accesos son idénticos (cada thread procesa
    un único elemento) y la coalescencia es perfecta en los dos casos.

    Para $m > 1$, el Kernel A genera un patrón de acceso con stride $m$, lo que incrementa el número
    de transacciones de memoria por warp de forma proporcional a $m$ (hasta el límite de 32). El
    Kernel B mantiene accesos coalescentes en todas sus iteraciones, por lo que su consumo de ancho
    de banda es óptimo independientemente de $m$. En consecuencia, se espera que el Kernel B
    \textbf{supere} al Kernel A en rendimiento para $m > 1$, con una brecha que se amplía a medida
    que $m$ crece.
\end{exercise}

\begin{exercise}{3}
    Análisis dinámico, es decir, haciendo pruebas. Será la forma de comprobar el rendimiento real de
    los kernels, y comparar con los resultados del apartado 2.

    Para cada kernel, se buscará la mejor configuración de threads y bloques de threads. Para ello
    se debe variar el número de threads total (variando $m$, con valores $1$, $64$ y $256$) y el
    número de threads por bloque ($16$, $32$, $64$, $128$ y $256$).

    En las pruebas se considerará lo siguiente:
    \begin{itemize}
        \item 10 tamaños de los vectores (número de componentes $n$): entre $2^{16}$ y $2^{25}$.
        \item 10 ejecuciones para cada prueba.
        \item Datos de tiempo total de ejecución y tiempo debido sólo al kernel.
    \end{itemize}

    Se deben generar gráficas (y tablas) con los datos de tiempo medido, en las que se comparen los
    resultados para los dos kernels, realizando un análisis de rendimiento a partir de los mismos.
    Además, se compararán esos tiempos con el tiempo del código secuencial.

    Para comprender los resultados que se obtienen, se deben observar también los accesos a memoria,
    para lo cual se usará la aplicación \textit{ncu} de NVIDIA.

    \begin{center}
        \noindent\rule{4cm}{0.2pt}
    \end{center}


\end{exercise}

\end{document}